\chapter{Inertial Sensors and Their Errors}
\label{ch:imu}

\begin{objectives}
\begin{itemize}
  \item Understand what MEMS accelerometers and gyroscopes physically measure, especially specific force.
  \item State the IMU measurement model of the project and the role of each term.
  \item Know the main IMU error sources and which ones the filter models.
  \item Convert between continuous-time noise densities, discrete per-sample noise, and datasheet units.
  \item Read an Allan deviation plot and use it to identify noise parameters.
\end{itemize}
\end{objectives}

% ------------------------------------------------------------------------------
\section{How MEMS inertial sensors work}

A MEMS (micro-electro-mechanical system) \textbf{accelerometer} is a tiny proof mass suspended on silicon
springs. When the case accelerates, the springs must deflect to drag the mass along. The deflection,
measured capacitively, is proportional to the force the springs exert per unit mass. A MEMS
\textbf{gyroscope} keeps a proof mass vibrating along one axis. When the sensor rotates, the Coriolis
force $-2m\,\omega\times\dot r$ pushes the mass along a perpendicular axis with an amplitude proportional to the
angular rate, and that secondary motion is measured. A six-axis IMU packages three of each, nominally
orthogonal, with electronics that filter, sample and timestamp the signals.

\section{Specific force: what an accelerometer measures}

Newton's second law for the proof mass $m$ in an inertial frame reads
\begin{equation}
  m\,a = F_{\text{spring}} + m\,g,
\end{equation}
where $a$ is the acceleration of the mass (equal to that of the case once transients settle) and $g$ is the
gravitational acceleration. The sensor measures the spring force, i.e.\ the \emph{specific force}
\begin{equation}
  f = \frac{F_{\text{spring}}}{m} = a - g .
\end{equation}
Gravity acts on the mass and the case equally and cannot be felt directly: this is Einstein's equivalence
principle. Two cases build intuition.
\begin{itemize}
  \item \textbf{At rest on a table}, $a = 0$, so $f = -g$. With $g_w = [0, 0, -9.81]\T$, the accelerometer reads
        $+9.81~\mathrm{m/s^2}$ \emph{upward}: it senses the table pushing it up.
  \item \textbf{In free fall}, $a = g$, so $f = 0$. A falling IMU reads zero, whatever its velocity.
\end{itemize}
To recover the kinematic acceleration from the measurement, the estimator rotates the specific force into
the world frame and \emph{adds gravity back}: $a_w = R_{wb}f_b + g_w$. A sign error here produces a vertical
acceleration error of $2g \approx 20~\mathrm{m/s^2}$ and an obviously diverging altitude, which is
fortunately easy to recognize.

% ------------------------------------------------------------------------------
\section{The measurement model}

\begin{convention}[IMU model, AGENTS.md C6]
\begin{equation}
  \omega_m = \omega_b + b_g + n_g, \qquad
  a_m = R_{wb}\T\big(a_w - g_w\big) + b_a + n_a, \qquad
  \dot b_g = n_{b_g}, \quad \dot b_a = n_{b_a},
  \label{eq:imu-model}
\end{equation}
with $\omega_m$, $a_m$ the gyroscope and accelerometer readings in the body frame, $\omega_b$ the true angular rate,
$a_w$ the true acceleration of the body in the world frame, $b_g$, $b_a$ the biases (body frame), and
$n_g, n_a, n_{b_g}, n_{b_a}$ zero-mean white noises with densities $\sigma_g, \sigma_a, \sigma_{b_g}, \sigma_{b_a}$. The noise
vector is ordered $w = [n_a, n_g, n_{b_a}, n_{b_g}] \in \R^{12}$.
\end{convention}

The model has two kinds of error. \emph{White noise} ($n_g$, $n_a$) is independent from sample to sample and
averages out. \emph{Bias} ($b_g$, $b_a$) is a slowly varying offset that does not average out and is
therefore estimated as part of the state. Its slow drift is modeled as a random walk driven by
$n_{b_g}$, $n_{b_a}$. Over the minutes of a EuRoC sequence the biases change little but measurably.

% ------------------------------------------------------------------------------
\section{Error sources}

\begin{table}[t]
\centering\small
\begin{tabularx}{\textwidth}{@{}l X l@{}}
\toprule
\textbf{Error} & \textbf{Description} & \textbf{In the model?} \\
\midrule
White noise & sample-to-sample electronic and thermo-mechanical noise (angle/velocity random walk) & yes, $n_g$, $n_a$ \\
Turn-on bias & constant offset that differs at each power-up & yes, initial $b$ and $P_0$ \\
In-run bias drift & slow wandering of the offset & yes, random walk \\
Bias instability & flicker ($1/f$) noise floor of the bias & approximated by the random walk \\
Scale factor & output $= (1 + s)\times$ input & no \\
Axis misalignment & sensing axes not exactly orthogonal / aligned with the case & no \\
$g$-sensitivity & gyro output depends on acceleration & no \\
Temperature effects & bias and scale factor change with temperature & no (absorbed by bias walk) \\
Quantization & finite resolution of the ADC & no (below white noise) \\
Saturation, vibration & clipping on impacts; aliasing of high-frequency vibration & no \\
Earth rotation & gyros sense $7.3\times10^{-5}$~rad/s & no (\cref{ch:rotations}) \\
\bottomrule
\end{tabularx}
\caption{IMU error sources and how the project's model treats them.}
\label{tab:imu-errors}
\end{table}

\Cref{tab:imu-errors} lists the main error sources. Unmodeled errors do not disappear: they appear as
apparent extra noise or apparent bias drift. A filter that ignores them must compensate with larger noise
parameters than the white-noise floor alone would suggest, which is why tuning is needed even with a
well-calibrated sensor. High-accuracy systems estimate scale factors and misalignments too, and EuRoC's
calibration with Kalibr~\cite{furgale2013kalibr} removes most static misalignment beforehand.

% ------------------------------------------------------------------------------
\section{Noise parameters and their units}

The noise densities are continuous-time quantities (\cref{sec:ct-noise}). For EuRoC's ADIS16448, the
\code{imu0/sensor.yaml} file gives the values in \cref{tab:euroc-imu}.

\begin{table}[h]
\centering\small
\begin{tabular}{@{}l l l l@{}}
\toprule
YAML key & Symbol & Value & Units \\
\midrule
\code{gyroscope\_noise\_density} & $\sigma_g$ & $1.6968\times10^{-4}$ & $\mathrm{rad/s/\sqrt{Hz}}$ \\
\code{accelerometer\_noise\_density} & $\sigma_a$ & $2.0\times10^{-3}$ & $\mathrm{m/s^2/\sqrt{Hz}}$ \\
\code{gyroscope\_random\_walk} & $\sigma_{b_g}$ & $1.9393\times10^{-5}$ & $\mathrm{rad/s^2/\sqrt{Hz}}$ \\
\code{accelerometer\_random\_walk} & $\sigma_{b_a}$ & $3.0\times10^{-3}$ & $\mathrm{m/s^3/\sqrt{Hz}}$ \\
\code{rate\_hz} & $1/\Dt$ & 200 & Hz \\
\bottomrule
\end{tabular}
\caption{EuRoC IMU noise parameters (\code{imu0/sensor.yaml}). Verify against your copy of the dataset.}
\label{tab:euroc-imu}
\end{table}

The units follow from the definitions. A gyro rate in rad/s with a white noise density has units
$\mathrm{rad/s}\cdot\sqrt{\mathrm{s}} = \mathrm{rad/s/\sqrt{Hz}}$. A gyro bias random walk is the integral of white
noise with units of $\mathrm{rad/s^2}$, so its density is in $\mathrm{rad/s^2/\sqrt{Hz}}$, and after $t$ seconds the
bias has standard deviation $\sigma_{b_g}\sqrt t$ in rad/s.

\paragraph{Discrete equivalents.} At sample period $\Dt$:
\begin{equation}
  \sigma_{g,d} = \frac{\sigma_g}{\sqrt\Dt}, \qquad \sigma_{a,d} = \frac{\sigma_a}{\sqrt\Dt}, \qquad
  \sigma_{b_g,d} = \sigma_{b_g}\sqrt\Dt, \qquad \sigma_{b_a,d} = \sigma_{b_a}\sqrt\Dt .
\end{equation}
For EuRoC: $\sigma_{g,d} = 2.4\times10^{-3}$~rad/s per sample, $\sigma_{a,d} = 2.8\times10^{-2}~\mathrm{m/s^2}$ per sample,
and the gyro bias diffuses by $1.4\times10^{-6}$~rad/s per sample. In the filter you never need these
directly, because $Q_d = G_cQ_cG_c\T\Dt$ uses the densities, but they are what a synthetic data generator must
add to each sample.

\paragraph{Datasheet units.} Manufacturers often quote angle random walk (ARW) in
$^\circ/\sqrt{\mathrm{h}}$ and velocity random walk (VRW) in $\mathrm{m/s/\sqrt{h}}$. Since
$1\ ^\circ/\sqrt{\mathrm{h}} = \frac{\pi/180}{60}~\mathrm{rad/\sqrt{s}} = 2.909\times10^{-4}~\mathrm{rad/s/\sqrt{Hz}}$, EuRoC's
$\sigma_g$ is an ARW of $0.58\ ^\circ/\sqrt{\mathrm{h}}$, and its $\sigma_a$ is a VRW of $0.12~\mathrm{m/s/\sqrt{h}}$.
Bias instability is quoted in $^\circ/\mathrm{h}$ and is \emph{not} the same thing as the random walk density.

% ------------------------------------------------------------------------------
\section{The Allan variance}
\label{sec:allan}

How are these parameters measured? Record the sensor at rest for a few hours, then compute the
\emph{Allan variance}. Split the record of rate samples into clusters of duration $\tau$, average each
cluster to get $\bar\omega_k(\tau)$, and compute half the mean squared difference of consecutive cluster
averages:
\begin{equation}
  \sigma_A^2(\tau) = \tfrac12\,\E\big[(\bar\omega_{k+1}(\tau) - \bar\omega_k(\tau))^2\big].
\end{equation}
Each noise process has a characteristic signature on a log-log plot of the Allan deviation
$\sigma_A(\tau)$:
\begin{center}\small
\begin{tabular}{@{}l l l@{}}
\toprule
Process & $\sigma_A^2(\tau)$ & slope of $\sigma_A$ \\
\midrule
quantization noise & $3Q^2/\tau^2$ & $-1$ \\
white noise (density $\sigma$) & $\sigma^2/\tau$ & $-\tfrac12$ \\
bias instability ($B$) & $\approx (0.664 B)^2$ & $0$ \\
bias random walk (density $\sigma_b$) & $\sigma_b^2\tau/3$ & $+\tfrac12$ \\
rate ramp ($R$) & $R^2\tau^2/2$ & $+1$ \\
\bottomrule
\end{tabular}
\end{center}
White noise is read off the $-\tfrac12$ line at $\tau = 1$~s, and the random walk density off the $+\tfrac12$ line
at $\tau = 3$~s. The minimum of the curve marks the averaging time beyond which a stationary sensor's
bias starts to wander, and its height is related to the bias instability.

\begin{figure}[t]
  \centering
  \includegraphics[width=\textwidth]{fig_allan.pdf}
  \caption{Allan deviation of a simulated 4-hour static recording with EuRoC's noise parameters. The
  $-\tfrac12$ slope of white noise dominates short averaging times and the $+\tfrac12$ slope of the bias random
  walk dominates long ones. Reading the asymptotes at $\tau = 1$~s and $\tau = 3$~s recovers the
  densities.}
  \label{fig:allan}
\end{figure}

\Cref{fig:allan} shows the procedure on simulated data. Real sensors add a flat bias-instability region
near the minimum, which the random-walk model only approximates. The IEEE standard
procedure~\cite{ieee952} and El-Sheimy et al.~\cite{elsheimy2008allan} give details.

% ------------------------------------------------------------------------------
\section{Tuning in practice}

The densities from an Allan analysis describe a stationary sensor in a quiet lab. On a flying vehicle,
vibration, temperature changes, scale-factor and misalignment errors, timing jitter and integration error
all add to the effective noise. Common practice is to start from the Allan values (or the dataset's values)
and inflate them by a modest factor until the filter is consistent, as judged by the statistical tests of
\cref{ch:consistency}. Two rules keep tuning honest.
\begin{enumerate}
  \item \textbf{Change physically meaningful parameters only}, and keep them within a plausible range of
        their measured values.
  \item \textbf{Judge on consistency statistics} (NEES on synthetic data, NIS and $3\sigma$ coverage on real data),
        not on trajectory error alone. Tuning for the lowest trajectory error on the test sequence is
        overfitting.
\end{enumerate}

\begin{pitfall}[mixing up the two discretizations]
The white-noise variance per sample is $\sigma^2/\Dt$ (larger at higher rate), while the random-walk variance per
step is $\sigma_b^2\Dt$ (smaller at higher rate). Swapping them is a common bug in synthetic data generators,
and it produces noise that is wrong by a factor of $1/\Dt^2 = 40\,000$ at 200~Hz.
\end{pitfall}

\section{The EuRoC IMU stream}

EuRoC's \code{imu0/data.csv} contains one row per sample: \code{timestamp [ns], w\_x, w\_y, w\_z [rad/s],
a\_x, a\_y, a\_z [m/s\textsuperscript{2}]}, i.e.\ $\omega_m$ and $a_m$ in $B$ (AGENTS.md C7). Samples arrive at
nominally 200~Hz, hardware-synchronized with the cameras. The filter must nevertheless take $\Dt$ from
consecutive timestamps (AGENTS.md C5): real data has jitter and occasional gaps, which Milestone 3 logs.
The ground-truth file starts later than the IMU file, so the filter is initialized at the first
ground-truth sample at or after the stream start.

\begin{projectlink}[Milestones 3 and 4]
M3 parses \code{sensor.yaml} into exactly the parameters of \cref{tab:euroc-imu} and builds the event stream.
M4's synthetic trajectory generator produces noise-free IMU data from \eqref{eq:imu-model}; adding
noise with the discrete equivalents above turns it into the Monte Carlo generator of M5.
\end{projectlink}

\begin{takeaways}
\begin{itemize}
  \item An accelerometer measures specific force $f = a - g$: it reads $+g$ upward at rest and zero in free fall.
  \item $\omega_m = \omega_b + b_g + n_g$ and $a_m = R_{wb}\T(a_w - g_w) + b_a + n_a$, with random-walk biases.
  \item Noise densities have units of $[u]/\sqrt{\mathrm{Hz}}$; per-sample white noise is $\sigma/\sqrt\Dt$, per-step bias
        diffusion is $\sigma_b\sqrt\Dt$.
  \item The Allan deviation identifies white noise (slope $-\tfrac12$, read at 1~s) and bias random walk (slope
        $+\tfrac12$, read at 3~s).
  \item Unmodeled errors force noise inflation; tune on consistency, not on accuracy.
\end{itemize}
\end{takeaways}

\begin{exercises}
  \item An IMU lies on a table tilted by $10^\circ$ about its $x$-axis. Compute the ideal accelerometer reading.
        Which component reveals the tilt, and with what sign?
  \item A gyroscope datasheet gives ARW $= 0.2\ ^\circ/\sqrt{\mathrm{h}}$. Convert it to
        $\mathrm{rad/s/\sqrt{Hz}}$ and to the per-sample standard deviation at 400~Hz.
  \item Simulate 2 hours of a gyro with EuRoC parameters, compute the overlapping Allan deviation, and recover
        $\sigma_g$ and $\sigma_{b_g}$ from the asymptotes. How long a recording do you need to estimate $\sigma_{b_g}$ to 10\%?
  \item Show that the Allan variance of white noise with density $\sigma$ is $\sigma^2/\tau$, and of a random walk with
        density $\sigma_b$ is $\sigma_b^2\tau/3$.
  \item A scale-factor error of $s = 10^{-3}$ affects a gyro on a vehicle spinning at 1~rad/s. What equivalent bias
        does it produce? Why can't a constant-bias model absorb it during maneuvers?
\end{exercises}
